\section*{Resumen}
\addcontentsline{toc}{section}{Resumen}

Los equipos de operaciones de seguridad dependen de inteligencia de
amenazas estructurada (IOCs, mapeos a MITRE ATT\&CK, artefactos listos
para STIX/MISP) para actuar r\'apidamente sobre los hallazgos de
malware. Los agregadores de reputaci\'on amplios como VirusTotal
entregan esa estructura a una escala enorme, pero no reconcilian sus
muchas se\~nales entre s\'i, de modo que producir un mapeo de
t\'ecnica defendible y ligado a evidencia para una muestra concreta
todav\'ia suele recaer en un analista de ingenier\'ia inversa que
traza el c\'odigo manualmente y redacta un informe. Esta tesis
presenta \textbf{MalWhere}, una pipeline modular que combina
ingenier\'ia inversa est\'atica y detonaci\'on din\'amica en sandbox,
reconcilia sus hallazgos en mapeos de t\'ecnicas MITRE ATT\&CK con
confianza cuantificada, y exporta el resultado como paquetes STIX~2.1
y eventos MISP en vivo.

La pipeline se valida frente a tres familias de malware con
arquitecturas distintas: un binario de ransomware (Akira); un
cargador multietapa que suelta un driver rootkit, un m\'odulo
antiantivirus y un cliente C2 de acceso remoto (RoningLoader); y un
infostealer .NET rico en funcionalidades (WhiteSnakeStealer). Cada una
cuenta con datos de referencia obtenidos mediante ingenier\'ia inversa
manual a nivel de funci\'on, verificada con Ghidra, no mediante un
escaneo superficial. Dos contribuciones metodol\'ogicas van m\'as all\'a
de una pipeline est\'atica+din\'amica convencional: un modelo de
reconciliaci\'on de confianza cruzada entre fuentes \emph{y} entre
niveles de t\'ecnica, que trata una t\'ecnica observada tanto por
evidencia est\'atica como din\'amica, o por una t\'ecnica padre y una
subt\'ecnica desde fuentes distintas, como una se\~nal m\'as fuerte que
cualquiera de las dos por separado; y un \emph{bucle de
reenv\'io} que analiza de forma independiente cada componente que una
muestra multietapa suelta o inyecta, sin atribuir err\'oneamente las
capacidades de una carga soltada a la puntuaci\'on de confianza del
binario padre.

La metodolog\'ia de evaluaci\'on va m\'as all\'a de la precisi\'on y la
exhaustividad en bruto: cada falso positivo encontrado durante la
validaci\'on se rastre\'o individualmente hasta su evidencia exacta
(una combinaci\'on de importaciones concreta, el campo de datos crudo
de una firma de sandbox, un patr\'on de cadena) y se contrast\'o con el
texto completo del informe manual correspondiente, no solo con su
tabla resumen. De los 15 falsos positivos encontrados de esta forma en
las tres muestras, doce eran errores gen\'uinos de la pipeline con
causa ra\'iz identificada y corregida (diez sin m\'as, m\'as dos
correcciones de calibraci\'on de confianza cuya evidencia gen\'erica
subyacente permanece presente, honestamente rebajada, en lugar de
desaparecer), uno era un hueco en c\'omo se extrajeron los datos de
referencia y no un error de la pipeline, y dos se dejaron abiertos
deliberadamente a la espera de confirmaci\'on independiente en lugar
de resolverse solo con la evidencia propia de la pipeline. La
cobertura de detecci\'on est\'atica se ampli\'o posteriormente de unas
30 a aproximadamente 85 identificadores de t\'ecnica MITRE ATT\&CK,
cada adici\'on obtenida de las propias descripciones de t\'ecnica de
MITRE y verificada para no introducir nuevos hallazgos en las muestras
validadas antes de confiar en ella.

La pipeline resultante alcanza puntuaciones F1 a nivel de familia de
0.93 (Akira), 0.81 (WhiteSnakeStealer) y 0.89 (RoningLoader, contando
sus componentes reenviados), con precisi\'on alta en las
tres muestras. La tesis argumenta que este rastro de auditor\'ia, cada
regla ligada a evidencia verificable, cada error diagnosticado a una
causa concreta, cada extensi\'on comprobada frente a regresiones, es lo
que hace que el resultado de la pipeline sea utilizable precisamente
para la decisi\'on que las se\~nales no reconciliadas de un agregador de
reputaci\'on no pueden sostener por s\'i solas, y es la propiedad con
mayor probabilidad de transferirse a familias de malware que la
pipeline todav\'ia no ha visto.

\bigskip
\noindent\textbf{Palabras clave:} an\'alisis de malware, inteligencia
de amenazas, MITRE ATT\&CK, STIX, MISP, an\'alisis est\'atico,
an\'alisis din\'amico, detonaci\'on en sandbox, ingenier\'ia inversa,
reconciliaci\'on de confianza.
